% исправления 27.05.2005
\section{Нестандартный анализ}

Один из создателей теории моделей, А.\,Робинсон,%
\glossary{\Робинсон}
заметил, что с
её помощью можно придать точный смысл понятиям \лк бесконечно
малых\пк\ и \лк бесконечно больших\пк\ величин, с которыми
оперировали ещё Ньютон%
\glossary{\Ньютон}
и Лейбниц%
\glossary{\Лейбниц}
и которые затем были изгнаны и
заменены рассуждениями с эпсилонами и дельтами.

Это направление получило название \emph{нестандартного анализа}.%
\index{Нестандартный анализ|defin}
Целей тут две: во\д первых, упростить доказательства известных
теорем, во\д вторых, использовать методы нестандартного анализа
для получения новых результатов. Насколько эти цели достигнуты
за тридцать с лишним лет, прошедших с возникновения нестандартного
анализа?

Простота доказательств\т дело вкуса. Конечно, всякий
преподаватель курса математического анализа меч\-тает избавиться
от утомительных рассуждений с выбором достаточно малых
эпсилонов. Но если вместо этого нужно постоянно переходить от
модели к её элементарному расширению и обратно, лекарство может
оказаться страшнее болезни. Во всяком случае, \лк
нестандартные\пк\ учебники математического анализа для
нематематиков (один из них написан
Кейслером%
\glossary{\Кейслер}%
~\cite{keisler-nonstandard}) большого распространения не
получили.

Новые результаты действительно были получены; отметим, что
многие из них (но не все) впоследствии были передоказаны \лк
стандартными\пк\ методами, так что и здесь революции не
произошло.

Так или иначе, нестандартный анализ\т интересное приложение
теории моделей, и мы разберём несколько простых примеров. Более
подробно об этом можно прочесть в книгах
Дэвиса%
\glossary{\Дэвис}%
~\cite{davis-nonstandard} и
Успенского%
\glossary{\Успенский}%
~\cite{uspensky-nonstandard}, а также в последней
главе книги Робинсона%
\glossary{\Робинсон}%
~\cite{robinson}.

Идея нестандартного анализа проста. Среди действительных чисел,
увы, нет бесконечно малых (которые были бы меньше $1/n$ при всех
$n=1,2,3,\dots$)\т как говорят, поле вещественных чисел
удовлетворяет \emph{аксиоме Ар\-хи\-меда}.%
\index{Аксиома!Архимеда|defin}%
\glossary{\Архимед}
(Оригинальная
формулировка этой аксиомы: каковы бы ни были два отрезка, можно
отложить меньший из них столько раз, чтобы превзойти больший.)
Но можно рассмотреть элементарное расширение поля~$\mathbb{R}$,
в котором такие бесконечно малые элементы есть, и использовать
их для определения пределов, производных и прочего в исходном
поле.

Перейдём к формальным определениям. Мы будет рассматривать
вещественную прямую как модель очень богатой сигнатуры. Для
каждого отношения на $\bbR$ (c~произвольным числом аргументов)
введём свой предикатный символ. Получится $2^{\mathfrak{c}}$
предикатных символов. Кроме того, для каждой функции из $\bbR^n$
в $\bbR$ (при всех $n=0,1,2,\dots$) введём свой функциональный
символ. Это даст ещё $2^{\mathfrak{c}}$ символов.

Пусть $\nstR$\т любая нормальная интерпретация этой сигнатуры,
элементарно эквивалентная%
\index{Элементарная эквивалентность}
$\bbR$. Её можно считать полем,
расширяющим поле~$\bbR$. В самом деле, среди функциональных
символов есть двуместные символы для сложения и умножения. Они
задают некоторые операции в $\nstR$ и относительно этих операций
множество~$\nstR$ будет полем, так как аксиомы поля можно
записать в виде формул (эти формулы истинны в~$\bbR$, а потому и
в~$\nstR$). Аналогичное рассуждение с предикатом \лк меньше\пк\
показывает, что $\nstR$ является упорядоченным полем.

Это поле можно считать расширением поля~$\bbR$. В самом деле,
для каждого действительного числа $x$ в сигнатуре имеется
константа. Значения таких констант образуют подполе в $\nstR$,
изоморфное~$\bbR$. В самом деле, утверждения вида $a\hm\ne b$,
$a+b\hm=c$ и $ab\hm=c$ являются формулами, и переносятся
из~$\bbR$ в~$\nstR$. Аналогичным образом это вложение сохраняет
порядок.

Если поле $\nstR$ исчерпывается значениями констант из~$\bbR$,
то ничего интересного не получается. Поэтому мы будем
предполагать, что это не так. Возможность построить~$\nstR$, не
совпадающее с~$\bbR$, следует (например) из теоремы Лёвенгейма\ч
Сколема о повышении мощности.%
\index{Теорема!Левенгейма\ч Сколема о повышении мощности}
Другой способ: добавим в сигнатуру
новую константу $c$ и рассмотрим теорию
         $$
\Th(\bbR)+\{c>\bar a\mid a\in\bbR\},
         $$
где $\Th(\bbR)$\т множество всех истинных в $\bbR$ формул нашей
сигнатуры, а $\bar a$\т константа для числа $a$. Совместность
этой теории следует из теоремы компактности.%
\index{Теорема!компактности}
Любая
её модель годится в качестве $\nstR$, поскольку значение
константы $c$ больше всех элементов из~$\bbR$.

\begin{problem}
        %
Проведите это рассуждение подробно.
        %
\end{problem}

В дальнейшем мы предполагаем, что выбрана и зафиксирована
некоторая интерпретация~$\nstR$, являющаяся элементарным
расширением~$\bbR$ и не совпадающая с~$\bbR$. Её элементы мы
называем \emph{гипердействительными}%
   \index{Гипердействительное число|defin}%
   \index{Число!гипердействительное|defin}
числами.
Среди них есть и действительные числа, которые мы будем называть
также \emph{стандартными}%
   \index{Гипердействительное число!стандартное|defin}%
   \index{Стандартное гипердействительное число|defin}
элементами~$\nstR$. Остальные элементы $\nstR$
будут \emph{нестандартными}%
   \index{Гипердействительное число!нестандартное|defin}%
   \index{Нестандартное гипердействительное число|defin}
гипердействительными числами. (По нашему предположению таковые
существуют.)

Утверждение об элементарной эквивалентности $\nstR$
и~$\bbR$ называют \emph{принципом переноса}:%
\index{Принцип переноса|defin}
он позволяет
перенести истинность формулы из~$\bbR$ в~$\nstR$ (или наоборот).

Возможность переноса не ограничивается алгебраическими
свойствами. Например, в нашей сигнатуре есть функция $\sin$. В
интерпретации $\nstR$ ей соответствует функция, которую можно
было бы назвать \лк гипердействительным синусом\пк. Эта функция
продолжает обычный синус (для стандартных аргументов), поскольку
утверждения вида $\sin a = b$ для конкретных стандартных~$a$
и~$b$ можно перенести в~$\nstR$. Более того, она обладает
обычными свойствами синуса: скажем, гипердействительный синус
любого гипердействительного числа не превосходит единицы (в
смысле порядка на~$\nstR$), поскольку формула $\forall x\,{\sin
x \le 1}$ выдерживает перенос. Аналогично можно поступать и с
предикатами: например, предикат \лк быть натуральным числом\пк\
задаёт в $\nstR$ некоторое подмножество, элементы которого
естественно назвать \emph{гипернатуральными}%
\index{Гипердействительное число!гипернатуральное|defin}%
\index{Гипернатуральное число|defin}
числами. Гипернатуральные
числа делятся на стандартные (соответствующие обычным натуральным
числам в~$\bbR$) и нестандартные. (Мы увидим, что нестандартные
числа обязательно найдутся.) Множество гипернатуральных чисел
обозначается $\nstbb{\bbN}$.

Аналогично определяется множество $\nstbb{\bbZ}$ \emph{гиперцелых
чисел}%
\index{Гипердействительное число!гиперцелое|defin}%
\index{Гиперцелое число|defin}
и вообще множество $\nst{M}$ для любого множества~$M$
действительных чисел. (Множеству $M$ соответствует одноместный
предикатный символ; $\nst{M}$\т интерпретация этого символа в
$\nstR$.) Множество $\nst{M}$ называют нестандартным
расширением~$M$. В нём содержатся те же стандартные числа, что и
в~$M$ (формулы вида $a\hm\in M$ для стандартных чисел~$a$
переносятся), и, возможно, некоторые нестандартные числа.

Принцип переноса гарантирует, что для конечного~$M$
нестандартных элементов в~$\nst{M}$ не появится. В самом деле,
пусть, скажем, в $M$ ровно три элемента~$a$, $b$ и~$c$. Тогда
формула
        $$
 \forall x\,(M(x) \leftrightarrow ((x=a)\lor(x=b)\lor(x=c))),
        $$
в которой $M(x)$\т предикат, соответствующий множеству~$M$,
истинна в $\bbR$. По принципу переноса она истинна и в $\nstR$,
так что и $\nst{M}$ состоит из трёх элементов, являющихся
значениями констант~$a$, $b$ и~$c$ (отождествлённых со
стандартными действительными числами).

Впоследствии мы увидим, что бесконечное множество~$M$
обязательно приобретёт новые нестандартные элементы при переходе от $\bbR$
к~$\nstR$.

Несколько простых следствий принципа переноса:
        %
\begin{itemize}
        \item
$M\subset N\Leftrightarrow \nst{M}\subset\nst{N}$
(применяем принцип переноса к формуле $\forall x\,(M(x)\hm\to N(x))$);
        \item
$\nst{(M\cup N)}=\nst{M}\cup \nst{N}$,
(применяем принцип переноса к формуле
$\forall x((M(x)\lor N(x))\leftrightarrow K(x))$, где
$K$\т объединение $M$ и $N$);
        \item
аналогичные утверждения верны и для пересечения и разности множеств.
        %
\end{itemize}

\begin{problem}
        %
Покажите, что для счётного объединения аналогичное утверждение
может не быть верным и $\nst{(M_0\cup M_1\cup M_2\cup\dots)}$
может отличаться от
$(\nst{M_0}\cup\nst{M_1}\cup\nst{M_2}\cup\dots)$.
        %
\end{problem}

Нестандартные аналоги имеют не только множества, но и функции.
Мы уже говорили о нестандартном аналоге синуса. Точно так же
можно определить нестандартный аналог любой всюду определённой
функции (любого числа аргументов). Для не всюду определённых
функций (например, для функции квадратного корня) надо
рассмотреть её график как предикат (для корня это будет предикат
двух аргументов) и взять его нестандартный аналог. Этот
нестандартный аналог будет графиком частичной функции (ибо
свойство \лк быть графиком частичной функции\пк\ записывается
формулой). Соответствующая функция и будет нестандартным
аналогом исходной.

\begin{problem}
        %
Покажите, что (построенный по этой схеме) нестандартный
квадратный корень имеет областью определения множество
неотрицательных гипердействительных чисел и что
$(\sqrt{x})^2\hm=x$ для любого неотрицательного
гипердействительного числа~$x$.
        %
\end{problem}

\begin{problem}
        %
Покажите, что для всюду определённой функции два способа
её продолжения (как функции и через график) дают одну и
ту же функцию.
        %
\end{problem}

\begin{problem}
        %
Покажите, что если множество $A$ является областью определения
частичной функции $\varphi$, то его нестандартный аналог
$\nst{A}$ совпадает с областью определения
функции~$\nst{\varphi}$.
        %
\end{problem}

Мы будем часто опускать звёздочки в записях вида~$\nst{f}(x)$,
считая, что если речь идёт о значении функции
на гипердействительном числе, то подразумевается нестандартный
аналог этой функции. (Путаницы не будет, так как на стандартных
числах значения функции и её гипердействительного аналога
совпадают.)

\begin{problem}
        %
Абсолютную величину%
\index{Гипердействительное число!абсолютная величина|defin}
гипердействительного числа~$x$ можно определить
как $x$ при $x\hm\ge0$  и как $(-x)$ при $x\hm\le 0$. С~другой
стороны, можно рассмотреть нестандартный аналог
функции $x\hm\mapsto |x|$. Покажите, что получится одно и то же.
        %
\end{problem}

\begin{problem}
        %
Покажите, что поле гипердействительных чисел является
вещественно замкнутым%
\index{Поле!вещественно замкнутое}%
\index{Вещественно замкнутое поле}
(согласно определению на
с.~\pageref{real-closed-fields})
        %
\end{problem}

Любое гипердействительное число~$x$ можно представить в виде
суммы гиперцелого числа~$n$ и некоторого гипердействительного
числа~$\alpha$, для которого $0\le\alpha<1$. Чтобы
убедиться в этом, достаточно рассмотреть нестандартные аналоги
функций целой и дробной части. Принцип переноса гарантирует, что
они сохранят свои свойства. В частности, в сумме они дают
исходное число, а дробная часть всегда не меньше нуля и меньше
единицы.

Целью расширения было получить возможность рассматривать
бесконечно большие и бесконечно малые числа. Дадим
соответствующие определения.

Гипердействительное число $\alpha$, большее всех стандартных
чисел, называется \emph{положительным бесконечно большим} числом.%
\index{Гипердействительное число!положительное бесконечно
большое|defin}%
\index{Положительное бесконечно большое число|defin}
Аналогично определяются отрицательные бесконечно большие числа.

\begin{problem}
        %
Докажите, что
число $\alpha$ является отрицательным бесконечно большим
тогда и только тогда, когда $-\alpha$ является положительным
бесконечно большим.
        %
Докажите, что  $|\alpha|$ является положительным бесконечно
большим тогда и только тогда, когда $\alpha$ является либо
положительным, либо отрицательным бесконечно большим.
        %
\end{problem}

Гипердействительные числа, не являющиеся бесконечно большими,
называют конечными. Другими словами, гипердействительное
число~$x$ \emph{конечно},%
\index{Гипердействительное число!конечное|defin}%
\index{Конечное гипердействительное число|defin}
если оно находится в
промежутке $a\hm\le x\hm\le b$ со стандартными концами $a$ и
$b$.

Наконец, гипердействительное число называется \emph{бесконечно
малым},%
\index{Гипердействительное число!бесконечно малое|defin}%
\index{Бесконечно малое гипердействительное число|defin}
если его абсолютная величина меньше любого стандартного
положительного числа. (Согласно этому определению нуль тоже
является бесконечно малым числом.) Легко проверить, что
ненулевое число $e$ является бесконечно малым тогда и только
тогда, когда~$1/e$~бесконечно велико. В самом деле, пусть,
например, $e>0$ бесконечно мало. Тогда $1/e$ больше любого
стандартного числа $c>0$, так как $e\hm<1/c$. Остальные случаи
разбираются аналогично.

Сумма и произведение двух конечных чисел конечны. Если $\alpha$
по модулю меньше стандартного числа $a$, а $\beta$\т
стандартного числа~$b$, то $\alpha+\beta$ по модулю меньше
стандартного числа $a+b$, а $\alpha\beta$ по модулю меньше~$ab$.
(Неравенства в гипердействительных числах можно складывать и
умножать, так как обычные свойства неравенств записываются
формулами и допускают перенос.)

В обычном курсе математического анализа аналогом этого
рассуждения является утверждение о том, что сумма и произведение
ограниченных последовательностей ограничены. Другое стандартное
утверждение из курса анализа\т о произведении ограниченных и
бесконечно малых (сходящихся к нулю) последовательностей\т также
имеет естественный аналог: произведение конечного и бесконечно
малого гипердействительных чисел является бесконечно малым
гипердействительным числом. Доказательство также вполне
традиционно: если $\alpha$ не превосходит стандартного
числа~$a$, а $|\beta|$ меньше любого стандартного положительного
числа, то $\alpha\beta$ меньше любого стандартного
положительного~$e$, так как $|\beta|\hm<e/a$.

Два гипердействительных числа $\alpha,\beta$ 
\emph{бесконечно близки},%
\index{Гипердействительные числа!бесконечно близкие|defin}%
\index{Бесконечно близкие гипердействительные числа|defin}
если их разность бесконечно мала.
Обозначение: $\alpha\approx\beta$.

\begin{problem}
        %
Докажите, что если $\alpha\approx\beta$, то
$\alpha+\gamma\hm\approx\beta+\gamma$ для любого
гипердействительного~$\gamma$, а
$\alpha\gamma\hm\approx\beta\gamma$ для любого конечного
гипердействительного~$\gamma$. Покажите, что условие конечности
существенно.
        %
\end{problem}

\begin{problem}
        %
Покажите, что два конечных гипердействительных числа
бесконечно близки тогда и только тогда, когда между
ними нельзя вставить двух разных стандартных чисел.
        %
\end{problem}

Легко проверить, что отношение бесконечной близости является
отношением эквивалентности на множестве гипердействительных
чисел. Классы эквивалентности этого отношения иногда называют
\emph{монадами}%
\index{Монада|defin}
(термин, использовавшийся ещё Лейбницем).%
\glossary{\Лейбниц}

\begin{theorem}\label{standard-part}
        %
Всякое конечное гипердействительное число
бесконечно близко к некоторому стандартному числу.%
%
\end{theorem}

(Заметим, что обратное утверждение очевидно: всякое
гипердействительное число, бесконечно близкое к некоторому
стандартному $a$, конечно, поскольку содержится между
стандартными числами $a-1$ и $a+1$.)

\begin{proof}
        %
Пусть $\alpha$\т конечное гипердействительное число. Рассмотрим
множество множество $L$ всех стандартных действительных чисел,
меньших или равных $\alpha$, а также множество $R$ всех
стандартных действительных чисел, больших или равных $\alpha$.
Конечность числа~$\alpha$ гарантирует, что оба этих множества
непусты (если бы, скажем,~$R$~было пусто, то~$\alpha$ было бы
положительным бесконечно большим). Заметим, что $L$ и $R$ не
пересекаются (если только~$\alpha$ само не является стандартным,
и тогда доказывать нечего) и в объединении дают всё множество~$\bbR$.

По аксиоме полноты существует действительное число $a$, для
которого $L\hm\le a \hm\le R$. Покажем, что $\alpha-a$
бесконечно мало. Проверим, например, что для любого стандартного
$e\hm>0$ выполнено неравенство $\alpha-a\hm<e$, то есть
$\alpha<a+e$. Это понятно: если $a+e\le\alpha$, то $a+e\in
L$, что противоречит свойству $L\hm\le a$. По аналогичным
причинам $\alpha-a\hm>-e$.
        %
\end{proof}

Стандартное число~$a$, бесконечно близкое к конечному
гипердействительному~$\alpha$, называется \emph{стандартной
частью}%
\index{Гипердействительное число!стандартная часть|defin}%
\index{Стандартная часть гипердействительного числа|defin}
числа~$\alpha$. Стандартная часть определена однозначно,
так как два разных стандартных числа не могут быть бесконечно
близки к одному и тому же гипердействительному числу (тогда бы
они были близки друг к другу, что невозможно). Поэтому можно
ввести обозначение $\st\alpha$ для стандартной части конечного
числа~$\alpha$.

\begin{problem}
        %
Докажите, что если $\alpha$ и $\beta$ конечны, причём
$\alpha\hm\le\beta$, то и  $\st\alpha\hm\le\st\beta$.
        %
\end{problem}

\begin{theorem}
        %
Среди гипердействительных чисел есть ненулевые бесконечно малые,
а также бесконечно большие числа.
        %
\end{theorem}

\begin{proof}
        %
Напомним, что по нашему предположению $\nstR$ не совпадает с
$\bbR$, то есть существует некоторое нестандартное
гипердействительное число $\alpha$. Если $\alpha$ бесконечно,
то~$1/\alpha$\т искомое ненулевое бесконечно малое число. Если
$\alpha$ конечно, то $\alpha-\st(\alpha)$ будет искомым ненулевым
бесконечно малым числом (а обратное к нему будет бесконечно
большим).
        %
\end{proof}

Заметим, что при построении гипердействительных чисел с помощью
формул $c>a$ (для новой константы~$c$ и~всех стандартных~$a$) и
теоремы компактности существование бесконечно больших элементов
очевидно: таковым будет значение этой самой константы~$c$.

Теперь обратимся к натуральным и целым числам.

\begin{theorem}
        %
Существуют нестандартные гипернатуральные числа, при этом
все они бесконечно велики.
        %
\end{theorem}

(Таким образом, для гипернатуральных чисел конечность и
стандартность равносильны.)

\begin{proof}
        %
Всякое положительное действительное число есть сумма
натурального и числа из $[0,1)$. Принцип переноса гарантирует,
что всякое положительное гипердействительное число $\alpha$ есть
сумма гипернатурального $\nu$ и гипердействительного $\tau$, для
которого $0\hm\le\tau\hm< 1$. Возьмём~$\alpha$ бесконечно
большим, тогда и $\nu$ будет бесконечно большим. Первое утверждение
доказано.

Пусть теперь $\nu$\т конечное гипернатуральное число. По определению
конечности оно меньше некоторого стандартного числа~$a$,
скажем, числа~$5$. Но в стандартной модели верна формула
        \begin{multline*}
\forall x\, ( ((x\in\bbN) \land (x<5)) \to{} \\
   {}\to((x=0)\lor (x=1)\lor (x=2)\lor (x=3)\lor (x=4))).
        \end{multline*}
По принципу переноса она верна и в $\nstR$, поэтому число
$\nu$ совпадает с одним из стандартных чисел $0,1,2,3,4$.
        %
\end{proof}

\begin{problem}
        %
Покажите, что для всякого гипердействительного числа существует
большее его гипернатуральное.
        %
\end{problem}

\begin{problem}
        %
Рассмотрим гипернатуральные числа как упорядоченное мно\-же\-с\-т\-во.
Покажите, что оно изоморфно $\bbN + \bbZ\times F$, где
$F$\т плотное линейно упорядоченное множество без первого
и последнего элементов.%
\index{Плотное линейно упорядоченное множество без первого
и последнего элементов}
(Порядок на $\bbZ\times F$ такой: сравниваются
сначала вторые элементы, а при равенстве\т первые.)
        %
\end{problem}

Гипернатуральные числа позволяют говорить о бесконечно далёких
членах (стандартных) последовательностей действительных чисел.
Пусть $a_0,a_1,\dots$\т такая последовательность. Рассмотрим её
график, то есть множество пар $\langle 0,a_0\rangle,\langle
1,a_1\rangle,\dots$, как двуместный предикат. Утверждение о том,
что этот предикат задаёт график функции, определённой на
натуральных числах, можно записать в виде формулы. Принцип
переноса гарантирует, что гипердействительный аналог этого
предиката будет функцией, определённой на гипернатуральных
числах и принимающей гипердействительные значения. Значение этой
функции на гипернатуральном числе $n$ можно обозначать $a_n$, не
опасаясь путаницы (при стандартных~$n$ мы получаем одно и то
же).

Таким образом, любая последовательность приобретает\т помимо
своего желания\т бесконечный \лк хвост\пк.

\begin{problem}
        %
Покажите, что если две последовательности отличаются лишь
в конечном числе членов, то их бесконечные хвосты одинаковы.
        %
\end{problem}

Сейчас мы используем продолжение последовательностей для
доказательства такого факта:

\begin{theorem}
        %
Нестандартный аналог $\nst{A}$ множества~$A$ действительных чисел совпадает
с $A$ тогда и только тогда, когда множество~$A$ конечно.
        %
\end{theorem}

\begin{proof}
        %
Если~$A$ конечно, и, скажем, состоит из трёх элементов $p,q,r$,
то можно записать формулу
        $$
\forall x\, ((x\in A) \leftrightarrow ((x=p)\lor (x=q)\lor (x=r)).
        $$
По принципу переноса эта формула остаётся истинной в~$\nstR$,
так что~$\nst{A}$ состоит из тех же трёх элементов.

Пусть теперь~$A$ бесконечно. Покажем, что $\nst{A}$ содержит
элементы, не входящие в~$A$. Пусть $a_0,a_1,\ldots$\т
последовательность различных элементов множества~$A$. Напишем
формулу, которая утверждает, что все элементы этой
последовательности различны и принадлежат~$A$. По принципу
переноса все бесконечные члены этой последовательности (точнее,
её гипердействительного аналога) также различны,
принадлежат~$\nst{A}$ и отличаются от всех конечных членов
последовательности. Они и будут искомыми нестандартными
элементами~$\nst{A}$. В самом деле, бесконечный член $a_\nu$ при
бесконечном гипернатуральном $\nu$ не может совпасть с конечными
членами, а также не может совпасть со стандартным элементом
$a\hm\in A$, не входящим в исходную последовательность (ибо
утверждение \лк $a_n\ne a$ при всех $n$\пк\ записывается
формулой).
        %
\end{proof}

\emph{Галактикой}%
\index{Галактика гипердействительного числа|defin}%
\index{Гипердействительное число!галактика|defin}
гипердействительного числа $\alpha$ называют
множество всех гипердействительных $\beta$, для которых разность
$\alpha-\beta$ конечна.

\begin{problem}
        %
Покажите, что множество гипердействительных чисел разбивается на
галактики. Определите на галактиках естественное отношение линейного
порядка и покажите, что этот порядок плотный и не имеет наибольшего
и наименьшего элементов.
        %
\end{problem}

\begin{problem}
        %
Каждое действительное число $a$, не являющееся двоично\д
рациональным, можно единственным образом записать в виде
бесконечной двоичной дроби $\ldots,a_0a_1\dots$; другими
словами, ему соответствует последовательность нулей и единиц
(нас будет интересовать лишь дробная часть после запятой).
Фиксируем бесконечное гипернатуральное $\nu$ и рассмотрим те
числа $a$, у которых $a_\nu=0$. Покажите, что множество таких
чисел переходит в своё дополнение при симметрии относительно
любой двоично\д рациональной точки (другими словами, ${a\in
M}\hm\Leftrightarrow {r-a\notin M}$ для двоично\д рациональных $r$)
и потому не может быть измеримым по Лебегу.%
\glossary{\Лебег}
        %
\end{problem}

\begin{problem}
        %
Докажите, что гиперрациональными числами являются отношения
гиперцелых чисел и только они.
        %
Докажите, что каждое гипердействительное число бесконечно близко
к некоторому гиперрациональному числу.
        %
\end{problem}

Покажем теперь, как можно ввести
основные понятия математического анализа, используя бесконечно
малые и бесконечно большие числа.

\begin{theorem}
        %
Пусть $M\subset\bbR$. Множество $M$ ограничено (в обычном смысле)
тогда и только тогда, когда все элементы
его гипердействительного аналога конечны.
        %
\end{theorem}

Таким образом, в курсе нестандартного анализа можно определять
ограниченные множества как множества, не содержащие бесконечных
элементов.

\begin{proof}
        %
Если все элементы $M$ меньше некоторого стандартного $a$ по модулю,
то и все элементы $\nst{M}$ меньше того же~$a$ (принцип переноса),
поэтому в одну сторону утверждение очевидно.

Пусть теперь $M$ не ограничено (скажем, сверху). Тогда в $\bbR$
верно такое утверждение: для всякого $c$ найдётся элемент
множества~$M$, больший $c$. Применим принцип переноса и возьмём
бесконечно большое $c$. Получим, что в~$\nst{M}$ есть бесконечно
большой элемент.
        %
\end{proof}

\begin{problem}
        %
Покажите, что если все элементы множества $\nst{M}$ меньше некоторого
гипердействительного $c$, то~$M$ ограничено.
        %
\end{problem}

\begin{problem}
        %
Говорят, что множество $S$ гипердействительных чисел является
\emph{внутренним},%
\index{Множество!гипердействительных чисел!внутреннее|defin}%
\index{Внутреннее множество гипердействительных чисел|defin}
если оно есть гипердействительный аналог
некоторого множества $A$ действительных чисел. Покажите, что
множество конечных гипердействительных чисел не является
внутренним.
        %
\end{problem}

\begin{problem}
        %
Докажите, что множество $S\subset\nstR$ выразимо (в рассматриваемой
нами сигнатуре, содержащей символы для всех функций и предикатов
на множестве~$\bbR$) тогда и только тогда, когда оно является
внутренним.
        %
\end{problem}

Нестандартный анализ позволяет дать естественные определения
предельной точки и предела.

\begin{theorem}
        %
Число $a$ является предельной точкой%
\index{Предельная точка!последовательности|defin}%
\index{Последовательность!предельная точка|defin}
последовательности действительных чисел
$a_0,a_1,\dots$ тогда и только тогда, когда найдётся
бесконечно далёкий член последовательности, бесконечно
близкий к~$a$.
        %
\end{theorem}

(Бесконечно далёким членом последовательности мы называем значение $a_\nu$
при бесконечном гипернатуральном~$\nu$.)

\begin{proof}
        %
Если $a$ является предельной точкой, то для всякого положительного
$\varepsilon$ и всякого натурального $N$
найдётся натуральное $n\hm>N$, для которого $|a_n-a|\hm<\varepsilon$.
Применим принцип переноса, положив $\varepsilon$ бесконечно малым
и $N$ бесконечно большим. Получим искомый бесконечно близкий к~$a$
член с бесконечно большим гипернатуральным номером.

Напротив, если для некоторого натурального $N$ и для некоторого
$\varepsilon\hm>0$ все члены последовательности, начиная с $N$\д
го, отстоят от $a$ более чем на $\varepsilon$, то по принципу
переноса все бесконечно далёкие члены последовательности также
отстоят от $a$ более чем на~$\varepsilon$.
        %
\end{proof}

\begin{problem}
        %
Покажите, что число $a$ принадлежит замыканию множества%
\index{Замыкание множества|defin}%
\index{Множество!замыкание|defin}
$M\hm\subset\bbR$ тогда и только тогда, когда некоторый элемент
множества $\nst{M}$ бесконечно близок к~$a$.
        %
\end{problem}

\begin{problem}
        %
Как определить в терминах нестандартного анализа понятие
предельной точки множества%
\index{Предельная точка!множества|defin}%
\index{Множество!предельная точка|defin}
(в любой окрестности которой
бесконечно много членов множества)?
        %
\end{problem}

Теперь видно, что нестандартный анализ позволяет в два счёта
доказать теорему о том, что ограниченная
последовательность имеет предельную точку: в самом деле, любой
бесконечно далёкий член этой последовательности конечен, и его
стандартная часть будет предельной точкой!

\begin{theorem}
        %
Последовательность $a_0,a_1,\dots$ действительных чисел сходится%
\index{Сходящаяся последовательность|defin}%
\index{Последовательность!сходящаяся|defin}
к числу $a$ тогда и только тогда, когда все её бесконечно далёкие
члены бесконечно близки~к~$a$.
        %
\end{theorem}

\begin{proof}
        %
Пусть~$a$ является пределом. Тогда для всякого $\varepsilon$
найдётся $N$, начиная с которого все члены последовательности
отстоят от $a$ менее чем на~$\varepsilon$. В частности, все
бесконечно далёкие члены таковы и их расстояние до $a$ меньше
любого стандартного~$\varepsilon$.

Напротив, пусть $a$ не является пределом и для всякого~$N$
найдётся член $a_n$ с номером $n>N$, отстоящий от~$a$ более
чем на~$\varepsilon>0$ (пока что все параметры стандартны).
Применим принцип переноса, взяв $N$ бесконечно большим,
и найдём бесконечно далёкий член последовательности,
отстоящий от~$a$ более чем на стандартное~$\varepsilon>0$.
        %
\end{proof}

Приведём теперь нестандартные критерии стандартных
топологических понятий.

\begin{theorem}
        %
Множество $M\hm\subset\bbR$ открыто%
\index{Открытое множество|defin}%
\index{Множество!открытое|defin}
тогда и только тогда,
когда вместе со всякой точкой $m\hm\in M$ оно содержит
и всю её монаду, то есть все гипердействительные точки,
бесконечно близкие к~$m$.
        %
\end{theorem}

(Cтрого говоря, следовало бы сказать \лк его
нестандартный аналог~$\nst{M}$\пк\ вместо \лк оно\пк;
напомним также, что $M$ и $\nst{M}$ содержат одни и те же
стандартные числа.)

\begin{proof}
        %
Если $M$ открыто и содержит вместе с точкой $m\hm\in M$ её
$\varepsilon$\д ок\-рест\-ность, то монада  $m$ по принципу
переноса содержится в $M$.

Если же некоторая точка $m\hm\in M$ не является
внутренней и для всякого действительного $\varepsilon>0$
найдётся точка вне~$M$ на расстоянии меньше $\varepsilon$,
применим принцип переноса и возьмём бесконечно малое
$\varepsilon$. Мы получим число, бесконечно близкое к
$m$ и не лежащее в $\nst{M}$.
        %
\end{proof}

Переходя к дополнениям, получаем, что множество $M\subset\bbR$
замкнуто тогда и только тогда, когда любая стандартная точка,
бесконечно близкая к некоторой точке из~$\nst M$, принадлежит
$M$.

На прямой компактными будут замкнутые ограниченные множества.
Соединим нестандартные критерии замкнутости и ограниченности:

\begin{theorem}
        \label{nonstandard-compactness}%
Множество $M\subset \bbR$ компактно%
\index{Компактное множество|defin}%
\index{Множество!компактное|defin}
тогда и только тогда,
когда любой элемент множества~$\nst{M}$ бесконечно близок
к некоторому (стандартному) элементу множества~$M$.
        %
\end{theorem}

\begin{proof}
        %
В самом деле, ограниченность означает, что любой элемент множества~$M$
конечен, то есть бесконечно близок к стандартному числу, а замкнутость
позволяет заключить, что это число принадлежит~$M$.
        %
\end{proof}

\begin{problem}
        %
Используя полученные только что критерии, покажите,
что любой отрезок $[a,b]$ действительной прямой
компактен, а любой интервал~$(a,b)$ открыт.
        %
\end{problem}

\begin{problem}
        %
Покажите, используя нестандартный критерий открытости, что
объединение любого числа открытых множеств открыто.
(Напоминание: гипердействительный аналог объединения может не
совпадать с объединением гипердействительных аналогов!)
        %
\end{problem}

\begin{problem}
        %
Покажите, что пересечение двух (или любого конечного числа)
открытых множеств открыто. (Где используется конечность?)
        %
\end{problem}

\begin{problem}
        %
Докажите, что последовательность фундаментальна тогда и
только тогда, когда любые два её бесконечных члена бесконечно
близки (предварительно уточнив формулировку этого утверждения).
        %
\end{problem}

\begin{problem}
        %
Докажите, что всякая фундаментальная последовательность сходится,
используя приведённый критерий фундаментальности. (Указание.
Ограниченность приходится доказывать, исходя из стандартных
определений.)
        %
\end{problem}

\begin{problem}
        %
Докажите, что если последовательность ограничена и имеет
единственную предельную точку, то она сходится (к этой точке).
        %
\end{problem}

\begin{problem}
        %
Докажите, что ограниченная возрастающая последовательность имеет
предел.
        %
\end{problem}

\medskip

Перейдём к функциям действительного переменного и дадим
не\-стан\-дарт\-ное определение предела (аналогичное приведённому
выше для последовательностей).

\begin{theorem}
        %
Число $b\in\bbR$ есть предел функции%
\index{Предел функции|defin}%
\index{Функция!предел|defin}
$f\colon \bbR\hm\to \bbR$ в
точке $a\hm\in\bbR$ тогда и только тогда, когда $f(x)\hm\approx
b$ для всех $x$, бесконечно близких к~$a$, но
отличных~от~$a$.
        %
\end{theorem}

\begin{proof}
        %
Пусть функция $f$ имеет предел $b$ согласно $\varepsilon$\д
$\delta$\д определению и для всякого $\varepsilon>0$ найдётся
$\delta>0$ с нужными свойствами. Бесконечно близкое к~$a$
число~$x$ попадает в $\delta$\д окрестность точки~$a$ при любом
стандартном~$\delta\hm>0$, поэтому $f(x)$ попадает в
$\varepsilon$\д окрестность точки $b$.

Напротив, если при некотором $\varepsilon\hm>0$ для любого
$\delta>0$ найдётся точка $x$, для которой $|x-a|\hm<\delta$, но
$|f(x)-b|\hm\ge\varepsilon$, то можно применить принцип переноса
(для данного стандартного~$\varepsilon$) и взять бесконечно
малое $\delta$.
        %
\end{proof}

Непосредственным следствием является нестандартный критерий
непрерывности: функция $f\colon \bbR\to \bbR$ непрерывна%
\index{Непрерывная функция|defin}%
\index{Функция!непрерывная|defin}
в (стандартной) точке $a$ тогда и только тогда, когда $f(x)\approx
f(a)$ для всех $x$, бесконечно близких к $a$.

Для функции, определённой на некотором множестве
$M\hm\subset\bbR$, критерий непрерывности в точке $m\in M$
выглядит так: $f(m)\hm\approx f(m')$ для всякой точки
$m'\hm\in\nst{M}$, бесконечно близкой к $m$.

\begin{problem}
        %
Проверьте это.
        %
\end{problem}

Поучительно понять, чем это свойство отличается от равномерной
непрерывности.

\begin{theorem}
        %
Функция $f\colon\bbR\to\bbR$ равномерно непрерывна%
\index{Равномерно непрерывная функция|defin}%
\index{Функция!равномерно непрерывная|defin}
на множестве
$M\subset\bbR$ тогда и только тогда, когда для всех $x,y\hm\in
\nst M$ выполнено ${x\approx y}\hm\Rightarrow {f(x)\approx
f(y)}$.
        %
\end{theorem}

\begin{proof}
        %
Пусть выполнено обычное $\varepsilon$\д $\delta$\д определение
равномерной непрерывности. Бесконечно близкие точки $x,y$ отличаются
менее чем на (стандартное) $\delta$, а потому их образы
отличаются менее чем на $\varepsilon$. Это верно для любого
стандартного~$\varepsilon$, поэтому $f(x)\hm\approx f(y)$.

Обратно, если функция не является равномерно непрерывной, то для
некоторого $\varepsilon$ и для любого $\delta$ найдутся точки,
отстоящие менее чем на $\delta$, образы которых отстоят более
чем на $\varepsilon$. Остаётся применить при данном
$\varepsilon$ принцип переноса и взять бесконечно
малое~$\delta$.
        %
\end{proof}

Чем это отличается от непрерывности во всех точках
множества~$M$? Непрерывность во всех точках~$M$ означает, что
для любого стандартного $m\hm\in M$ и любого бесконечно близкого
к нему $m'\hm\in \nst{M}$ мы имеем $f(m)\approx f(m')$.
Отсюда следует, что для любых $m',m''\hm\in\nst{M}$, бесконечно
близких к некоторому стандартному $m\hm\in M$, выполнено
$f(m')\hm\approx f(m'')$. Но в множестве $\nst{M}$ могут быть
бесконечно близкие элементы, стандартная часть которых не лежит
в~$M$ (или вообще не имеющие стандартной части, то есть
бесконечные). Легко понять, что для компактного~$M$ такого быть
не может (стандартная часть любого элемента $m'\hm\in\nst{M}$
принадлежит $M$ согласно теореме~\ref{nonstandard-compactness}).
Тем самым мы получили (почти что тривиальное) нестандартное
доказательство классической теоремы: непрерывная на компакте
функция равномерно непрерывна.

Вот ещё несколько \лк нестандартных\пк\ доказательств
стандартных (во всех смыслах этого слова) теорем из курса
математического анализа.

\begin{theorem}
        %
Функция, непрерывная на отрезке и принимающая значения разных
знаков на его концах, имеет нуль на этом отрезке.
        %
\end{theorem}

\begin{proof}
        %
Разделим отрезок на $n$ равных частей. Среди них найдётся часть,
на которой функция меняет знак. По принципу переноса и при
делении отрезка на бесконечное гипернатуральное число частей
найдётся часть, на которой функция меняет знак. Но концы этой
части бесконечно близки к некоторой стандартной точке отрезка.
Эта точка будет нулём функции (если в ней функция, скажем,
положительна, то по непрерывности в бесконечно близких к ней
концах отрезка изменения знака функция будет положительной).
        %
\end{proof}

\begin{theorem}
        %
Непрерывная во всех точках компакта функция ограничена на нём.
        \index{Функция!непрерывная на компакте}%
        %
\end{theorem}

\begin{proof}
        %
Пусть функция~$f$ непрерывна на компакте~$M$. Следуя
нестандартному критерию ограниченности, мы должны показать, что
значения функции $\nst{f}$ во всех точках~$\nst{M}$ конечны. Но
всякая точка $x\hm\in\nst{M}$ бесконечно близка к некоторой
стандартной точке $y\hm\in M$ (компактность), а потому
$\nst{f}(x)\hm\approx\nst{f}(y)$ (непрерывность),
поэтому~$\nst{f}(x)$ конечно.
        %
\end{proof}

Обратите внимание, что мы пользовались аксиомой полноты (для
множества~$\bbR$) только один раз, при доказательстве
теоремы~\ref{standard-part}. Это и не удивительно, поскольку из
утверждения этой теоремы следует аксиома полноты.

\begin{problem}
        %
Убедитесь в этом, следуя такой схеме. Пусть $A$\т произвольное
ограниченное
множество действительных чисел. Покажите (стандартными
рассуждениями), что для любого $\varepsilon$
найдётся число $c$, являющееся верхней гранью, для которого
$c\hm-\varepsilon$ не будет верхней гранью. Примените
принцип переноса, взяв бесконечно малое~$\varepsilon$ и
рассмотрев стандартную часть соответствующего числа~$c$.
        %
\end{problem}

\begin{problem}\label{derivative-non-standard}
        %
Докажите, что производная%
\index{Производная функции|defin}%
\index{Функция!производная|defin}
стандартной функции
$f\colon\bbR\hm\to\bbR$ в стандартной точке~$a$ равна
стандартному числу~$b$ тогда и только тогда, когда
${(f(a+h)-f(a))/h}\hm\approx b$ для всех бесконечно малых ${h\ne
0}$.
        %
\end{problem}

\begin{problem}
        %
Покажите, что $(x^n)'=nx^{n-1}$ согласно нестандартному
определению производной (предыдущая задача).
        %
\end{problem}

\begin{problem}
        %
Как использовать нестандартный анализ для определения
понятия интеграла?%
\index{Интеграл|defin}%
\index{Функция!интеграл|defin}
        %
\end{problem}

В наших примерах все рассмотрения были ограничены множеством
гипердействительных чисел. Это ограничение кажется
существенным\т не вполне ясно, каким образом можно применить те
же методы к произвольному топологическому пространству (в
котором нет бесконечно больших чисел). Тем не менее это
возможно, и об этом можно прочесть в книгах
Дэвиса%
\glossary{\Дэвис}%
~\cite{davis-nonstandard} или
Ус\-пенского%
\glossary{\Успенский}%
~\cite{uspensky-nonstandard}.
